\documentclass[11pt,letterpaper]{article}
\usepackage[margin=.7in,top=.78in,bottom=.7in]{geometry}
\usepackage[T1]{fontenc}
\usepackage{lmodern,amsmath,amssymb,array,booktabs,fancyhdr,hyperref}
\hypersetup{hidelinks}
\renewcommand{\familydefault}{\sfdefault}
\pagestyle{fancy}\fancyhf{}
\fancyhead[L]{\small Week 75 / Twists on a cylinder / Adult guide}
\fancyfoot[L]{\small Bellingham Math Circle / Week 75 / GGT75-FAC-v1}
\fancyfoot[R]{\small\thepage}
\renewcommand{\headrulewidth}{0pt}\renewcommand{\footrulewidth}{0pt}
\setlength{\parindent}{0pt}\setlength{\parskip}{6pt}
\setlength{\headheight}{14pt}
\newcommand{\sect}[1]{\par\vspace{5pt}{\large\bfseries #1}\par}
\newcommand{\prob}[1]{\par\vspace{3pt}\textbf{Problem #1}\par}

\begin{document}
\sect{The mathematics first}
A cylinder surface is an annulus: $(\mathbb R/\mathbb Z)\times[0,1]$.
Horizontal coordinate is measured in turns, so adding an integer returns to the
same place. A route is a simple arc from the fixed mark A on the bottom rim to
the fixed mark B on the top, with all other points strictly between the rims.
The printed tasks use upward routes. Their winding is the signed number of seam
crossings. In repeated rectangles, a lift starting at $(1/2,0)$ ends at
$(1/2+n,1)$ exactly when the winding is the integer $n$.

Winding cannot change during a continuous slide on the surface with both
endpoints fixed. Upward routes with the same winding can be slid into each other.
A full twist is the whole-surface map
\[
 T_k(\theta,t)=(\theta+kt\pmod1,t),\qquad k\in\mathbb Z.
\]
It fixes both rims pointwise, adds $k$ to every joining route's winding, and has
inverse $T_{-k}$. Composition adds: $T_jT_k=T_{j+k}$. Thus a twist word's effect
is its number of pluses minus minuses; its shortest equivalent word has that
signed sum's absolute value as its length.

For two simple joining routes of windings $a,b$ with \emph{these same fixed
endpoints}, the smallest number of interior intersections is
\[
 i(a,b)=\max\{|a-b|-1,0\}.
\]
Exclude the shared endpoints; shared pieces and non-crossing touches are not
allowed final comparisons. Arbitrary drawings may have extra intersections.
A common twist preserves the minimum. Twisting only one route can increase,
decrease or leave it unchanged. Children may discover a recoverable winding
record, cancellation, and how a difference controls unavoidable crossings.

\textbf{Who this fits.} One Grades 4--5 prototype for children ready to track
signed crossings, transfer a route between a cylinder and repeated strips, and
keep endpoints fixed. An adult may read or record. Sliding and lift arguments
need a mathematically confident adult; the general lower bound is a later
extension. There is no separate younger packet or classroom-tested age-fit claim.

\sect{Materials and preparation}
At the older three-child table, use a pair and rotating checker with one adult.
Per pair: the four-page packet, one extra page 1 for cutting, two or three spare
copies of page 2, two contrasting yarns about $60$ inches long for later high
windings, removable tape, scissors, pencils and blank paper. Coil excess yarn;
shorter pieces suffice for initial small-winding play.

Print at actual size. The cut rectangle is $6$ by $2.4$ inches. Tape its left and
right edges together without overlap; the dotted quarter-height lines should
meet. Keep the cylinder still, with A and B vertically aligned away from the
seam. Secure yarn at the geometric marks while leaving spare length beyond both
ends. Rehearse feeding slack through the fastening without shifting either mark.
If tape grips too tightly or the sleeve buckles, improve the fastening or keep
a drawn route as the mathematical record. This physical operation is unrehearsed.

\sect{Launch and a feasible first visit}
Let children handle the tube and yarn. Trace the printed seam example: an exit
through the right edge re-enters the left edge at the \emph{same height}, counting
$+1$. Show a route with fixed A and B and let a partner check it. Demonstrate how
one continues into the next repeated rectangle before giving out page 2.
Spend the first visit on Problems 1--3: route building, recording and same-winding
sliding may use the full hour. A useful stopping point is two different-looking
routes sharing a winding and an explanation of why different windings cannot
slide together. Save twist words and intersection minimization for a return visit
or a pair already secure with the cylinder-to-strip bridge.
\newpage
\sect{Answers and optional hints}
\prob{1}
\textbf{All four requested windings exist.} In normalized repeated rectangles,
draw the straight lift $x(t)=1/2+nt$, $0\le t\le1$, for $n=0,1,-1,2$.
Project the route back to the cylinder by wrapping x modulo $1$. These are simple
upward routes with signed seam totals $0,+1,-1,+2$. To change appearance without
changing winding, add a small sideward bulge away from both endpoints.
For example, $x(t)=1/2+nt+\tfrac14t(1-t)$ has the same endpoints.

\textit{Hint:} count signed crossings as you travel from A to B, rather than
counting every contact with the seam as a positive turn. Move a route slightly
if it lies along the seam or merely touches it; the counting convention needs
clear transverse crossings. A drawing or recorded lift is more reliable than
remembering how the yarn was laid.

\prob{2}
\textbf{Winding $+2$ ends at B in copy 2; winding $-1$ at B in copy $-1$.}
With A at $(1/2,0)$, the endpoints are $(5/2,1)$ and $(-1/2,1)$.
For each winding, the straight and gently bulged lifts above give two examples.
Every passage through a vertical copy boundary corresponds to a seam crossing;
height stays unchanged when returning to the base rectangle.
The copy endpoint records the net winding even for routes with extra opposite
seam crossings. \textit{Hint:} ask the partner to point to the matching B copy
before smoothing a drawn route.

\prob{3}
\textbf{Two distinct routes with the same winding can slide together; a pair
with different windings cannot.} Use the straight and bulged winding-$+1$
routes as the first pair, and windings $0,+1$ as the second.
For upward lifts $x_0(t),x_1(t)$ with the same endpoint copy, the interpolation
$(1-s)x_0(t)+sx_1(t)$, $0\le s\le1$, is an explicit slide. Height t remains
unique on each route, so no route crosses itself; endpoints remain fixed.

During any continuous fixed-endpoint slide, the lifted top endpoint moves
continuously but must always be one of the separated B copies. It cannot jump
from one copy to another. This proves winding invariance and the impossibility
for different windings. A different-looking route is not necessarily a different
winding. \textit{Hint:} ``Which B copy would its endpoint have to reach?''

\prob{4}
\textbf{Words have the same effect exactly when their signed sums agree.}
From winding $0$, a word with p pluses and m minuses ends at winding $p-m$.
Its shortest replacement is $p-m$ pluses if positive, $m-p$ minuses if negative,
and the empty word if zero. Each twist changes winding by only $1$, so fewer
than $|p-m|$ twists cannot suffice; repeating the needed sign attains it.
For examples, $++-$ has winding $1$ and reduces to $+$;
$+--$ reduces to $-$; $+--+$ reduces to no twists.

The displayed shifts are fractions of a full turn: at heights
$0,1/4,1/2,3/4,1$, a plus twist shifts by those same turn amounts.
At the top a full turn fixes every point modulo $1$.
Adding shifts proves $T_jT_k=T_{j+k}$ on the whole surface, so cancellation
works on every route. Do not physically rotate a marked rim to imitate this map.

\prob{5}
\textbf{Any word with three pluses and three minuses works}, such as $+++---$.
It has net shift zero, so composition is $T_0$, the identity, regardless of the
starting winding. Conversely, a length-six word has net zero only if it contains
three of each sign. There are $\binom63=20$ such words if children choose to
catalog them. Trying each possible starting route is unnecessary because the
same signed sum is added to all of them.
\newpage
\prob{6}
\begin{center}
\begin{tabular}{@{}ccccccc@{}}\toprule
Windings & $(0,1)$ & $(0,2)$ & $(0,3)$ & $(-1,1)$ & $(-1,2)$ & $(2,4)$\\
Fewest interior crossings & 0 & 1 & 2 & 1 & 2 & 1\\\bottomrule
\end{tabular}
\end{center}
\textbf{Construction.} For unequal windings use straight lifts
$(1/2+at,t)$ and $(1/2+bt,t)$. Their projections meet at an interior height
exactly when $(a-b)t$ is an integer. For $d=|a-b|\ge1$, those heights are
$1/d,2/d,\ldots,(d-1)/d$, giving $d-1$ crossings. They are distinct interior
points and genuine crossings. When $a=b$, put one route slightly alongside the
other, with the sideward gap shrinking to zero only at A and B. There are then
zero interior crossings; laying identical yarns on top of each other is not a
valid zero-crossing answer.

\textbf{Why no drawing can do better.} For upward lifts, the horizontal
difference of the two routes varies continuously from $0$ to $b-a$. It must
attain each integer strictly between those values. At every such height the
projected routes meet, and different attained integers need different heights.
Thus at least $|a-b|-1$ interior intersections are forced if the difference has
absolute value at least $2$. For differences $0$ or $1$, zero is the lower bound.

The same bound holds even for proper simple arcs that wiggle up and down.
Cut along the first arc to form a rectangle, straighten its two cut edges, then
reglue. This is a boundary-fixing homeomorphism sending it to a vertical arc;
the other arc's relative winding becomes $b-a$. In repeated copies, the first
arc's lifts are integer-spaced vertical barriers. A lift of the second must cross
each barrier strictly between its starting and ending copies. Simplicity prevents
two such forced meetings from being the same projected interior point. Together
with the construction, this proves the minimum rather than just a pattern in
examples. This general cut-and-reglue argument can remain with the adult.

\textit{Hint:} compare the two top copies in a repeated strip, or first test
one nearly straight vertical route against a winding route. Do not announce
the difference rule before children have minimized several contrasting pairs.

\prob{7}
\textbf{The minima are 4 for $(2,7)$ and 0 for $(5,5)$.} The general rule is
$\max(|a-b|-1,0)$. Adding the same integer twist k replaces $(a,b)$ by
$(a+k,b+k)$, leaving the difference and the minimum unchanged.
A twist on only one route can increase, decrease or preserve the answer:
$(0,2)\to(0,3)$ changes $1\to2$;
$(0,2)\to(1,2)$ changes $1\to0$;
$(5,5)\to(6,5)$ stays $0\to0$.

\sect{Physical meaning and source limits}
The rigid tube is a stationary strand-routing model. Lifting yarn to construct
a new route is permitted; it is not a legal continuous surface slide and cannot
prove two surface routes disjoint. Visible yarn overlaps represent intersections
of their projected centerlines on the surface, not knot over/under information.
No literal twisting paper sleeve is required. Keep fixed endpoints, rim marks,
proper interiors and excluded shared endpoints throughout; changing them changes
the question and can change the formula.

Diana Davis, \emph{Billiards, Surfaces, and Geometry}, Chapter 4,
\href{https://dianadavis.github.io/billiards-book.pdf}{Problem 111 and Figure 83, printed p.\ 111 (PDF p.\ 129)}, supplies related torus-shear and reassembly imagery.
It does not give this fixed-endpoint annulus model or intersection formula;
those are independently developed and proved here.

\textbf{Review status:} matched to final four-page GGT75-S-v1 student pages.
The independent checker covers the exact printed instances, twist words and
rational composition tests; general results use the proofs above. The guide is
built, rendered and visually checked. Physical preparation, surface sliding,
yarn fit and classroom pacing remain unrehearsed and unpiloted.
\end{document}
